\chapter{Prioritized Planning and Space-Time \astar}
\label{ch:ch08}
% Function names for pseudocode (unique across the book: prefix Pp).
\SetKwFunction{PpPlan}{PrioritizedPlanning}
\SetKwFunction{PpStar}{CooperativeAStar}
\SetKwFunction{PpReserve}{Reserve}
\SetKwFunction{PpWindowed}{WindowedCooperativeAStar}

\begin{objectives}
\begin{itemize}
  \item Explain prioritized planning in one paragraph: agents are planned one after
        another, and each treats the paths of the agents before it as moving obstacles.
  \item Build a reservation table by hand, including the stay-at-target reservation, and
        run Cooperative \astar against it on a small grid.
  \item Show with small instances that the priority order matters, that no order may
        work even when a solution exists, and that the best order can still be suboptimal.
  \item State and prove the completeness result for well-formed instances.
  \item Implement prioritized planning for two to five agents, create deliberate conflict
        cases, and validate the resulting plans.
  \item Explain the hierarchical and windowed variants of Cooperative \astar, choose a
        priority order with a simple heuristic, and describe priority-based search.
\end{itemize}
\end{objectives}

% ---------------------------------------------------------------------
\section{Why this matters for a drone swarm}
\label{sec:ch08-motivation}

Picture a warehouse roof with forty inspection drones. Each drone has a start pad and an
inspection point, and the airspace between them is cut into a grid of cells by the
support pillars. \Cref{ch:ch07} taught you how to describe this situation as a
multi-agent path finding (MAPF)\index{MAPF} instance and how to detect the conflicts of
a plan. It also warned you that finding an optimal plan is NP-hard. Forty drones is far
beyond what a joint search over all agents can handle, and even the optimal solvers of
\cref{ch:ch09,ch:ch10} need time that grows quickly with the number of interacting
agents. The operator, however, wants the drones in the air within a second.

The oldest and simplest answer is \textbf{prioritized planning}\index{prioritized planning}:
line the drones up in some order, plan the first one as if it were alone, then plan the
second one while treating the first drone's path as a moving obstacle, then the third
one against the first two, and so on. Each planning step is a single-agent space-time
\astar search of the kind you implemented in \cref{ch:ch04}. Forty drones cost forty
searches. That is why prioritized planning, introduced by Erdmann and
Lozano-P\'erez~\cite{erdmann1987multiple} and made popular for grids by Silver's
\emph{Cooperative \astar}~\cite{silver2005cooperative}, is the workhorse of video-game
crowds, warehouse robots and large drone swarms. It is also the study plan's Week~3
exercise (\cref{ch:appA}): implement it for two to five agents and create deliberate
conflict cases. In the hybrid architecture of \cref{ch:ch24} the global layer is \cbs or
\ecbs; prioritized planning is the natural fallback when they run out of time, and its
reservation table is the bookkeeping that a replanning layer needs when one drone has to
re-route around the frozen plans of all the others.

The price of the speed is that the method is greedy. A drone that is planned early never
makes room for a drone that is planned late. This chapter shows exactly when that
matters: the order can change the cost, the method can fail although a solution exists,
and even its best order can be more expensive than the optimum. It also gives the
positive result that makes prioritized planning safe in practice: on \emph{well-formed}
instances, where every agent can reach its goal without touching another agent's start
or goal, prioritized planning succeeds in every order~\cite{cap2015prioritized}.

\paragraph{Roadmap.} \Cref{sec:ch08-intuition} explains the idea and the reservation
table. \Cref{sec:ch08-problem} fixes the notation and \cref{sec:ch08-algorithm} gives the
pseudocode of prioritized planning and of Cooperative \astar. \Cref{sec:ch08-example}
works through a three-agent instance in every one of its six orders.
\Cref{sec:ch08-properties} collects the properties: soundness, incompleteness,
suboptimality, complexity and the completeness theorem for well-formed instances.
\Cref{sec:ch08-variants} covers the hierarchical and windowed variants, priority
ordering heuristics with a small experiment, and priority-based search.
\Cref{sec:ch08-implementation,sec:ch08-drone} give the implementation notes and the
place of the method in the drone system.

% ---------------------------------------------------------------------
\section{The idea in plain words}
\label{sec:ch08-intuition}

Suppose you have already planned drone~1. Its path tells you where it will be at every
time step: cell $v_0$ at time $0$, cell $v_1$ at time $1$, and so on until it reaches
its goal, where it stays. Now you plan drone~2. Instead of searching the joint space of
both drones, you search only for drone~2, in space-time, and you forbid every
(cell, time) pair that drone~1 occupies. Drone~1 is a moving obstacle: a wall that
appears in one cell at one time step and moves on. If drone~2 wants to enter a cell
that drone~1 occupies at that moment, drone~2 waits a step or goes around. When drone~2
is done, both paths are frozen, and drone~3 plans against both. The bookkeeping that
stores the frozen paths is the \textbf{reservation table}\index{reservation table}: a
set of (cell, time) entries that are taken, plus a set of (edge, time) entries so that
two drones cannot swap cells by moving through each other.

\begin{figure}[tb]
  \centering
  \inputfigure{ch08/idea}
  \caption{Prioritized planning on a $6 \times 5$ grid with two pillars. Left: agent~1
  is planned first against an empty reservation table and takes a shortest path. Middle:
  the cells of that path are now reserved at the times agent~1 occupies them (shaded),
  and agent~2 plans around them. Right: agent~3 plans against both earlier paths; it
  waits one step before crossing the corridor between the pillars. Notice that the
  earlier agents never move for the later ones: waits and detours are always paid by
  whoever comes later in the order.}
  \label{fig:ch08-idea}
\end{figure}

\Cref{fig:ch08-idea} shows the three steps on a small map, and
\cref{fig:ch08-spacetime} in the worked example draws the table itself as a space-time
diagram. Three details matter and are easy to get wrong.
\begin{itemize}
  \item \textbf{Stay at target.}\index{stay-at-target} Under the stay-at-target
        convention of \cref{ch:ch07}, an agent that arrives at its goal at time $T$
        remains there for ever. Its goal cell must therefore be reserved for \emph{all}
        times $t \ge T$, not only for $t = T$. Forgetting this lets a later agent drive
        through a parked drone.
  \item \textbf{Edge reservations.} Two agents that exchange cells between $t$ and $t+1$
        never occupy the same cell at the same time, so vertex reservations alone do not
        catch this swapping conflict. When agent~1 moves from $u$ to $v$ between $t$ and
        $t+1$, the table also reserves the reverse edge $(v, u)$ at time $t$.
  \item \textbf{The goal of a later agent.} A later agent may only finish its path at a
        time $T$ after which its goal is never reserved. If a higher-priority agent
        passes through that goal cell at some time $t' > T$, the later agent must wait
        somewhere and arrive after $t'$.
\end{itemize}

\begin{keyidea}
Prioritized planning turns a $k$-agent problem into $k$ single-agent space-time searches.
Agent $i$ sees agents $1, \dots, i-1$ as moving obstacles stored in a reservation table,
and never sees agents $i+1, \dots, k$ at all. This is what makes it fast, and it is also
the only reason it can fail.
\end{keyidea}

Cooperative \astar\index{Cooperative A*}~\cite{silver2005cooperative} is the name Silver
gave to the combination of space-time \astar with a reservation table. The word
\emph{cooperative} is slightly optimistic: only the later agents cooperate, by yielding
to the earlier ones. Nevertheless, in open maps with few agents this one-sided
cooperation is all that is needed, and the plans it produces are usually close to
optimal. The trouble starts in corridors and near goals, as the counterexamples of
\cref{sec:ch08-properties} show.

% ---------------------------------------------------------------------
\section{Problem statement and notation}
\label{sec:ch08-problem}

We use the classical MAPF setting of \cref{ch:ch07} and of Stern et
al.~\cite{stern2019mapf}: an undirected graph $G = (V, E)$, usually a 4-connected grid,
and $k$ agents $a_1, \dots, a_k$ with distinct start vertices $s_i$ and distinct goal
vertices $g_i$. Time is discrete. In one time step an agent either moves along an edge
or waits in its cell; both actions cost one. A \textbf{path}\index{path!time-indexed}
for agent $i$ is a sequence $\pi_i = (\pi_i(0), \pi_i(1), \dots, \pi_i(T_i))$ with
$\pi_i(0) = s_i$, $\pi_i(T_i) = g_i$ and consecutive vertices equal or adjacent. The
agent stays at $g_i$ afterwards, so we define $\pi_i(t) = g_i$ for all $t > T_i$. The
cost of the path is its arrival time $T_i$ and the cost of a plan is the sum of costs
$\sumcost = \sum_i T_i$\index{sum of costs}; its makespan is $\makespan = \max_i T_i$.
Two paths $\pi_i$ and $\pi_j$ have a \textbf{vertex conflict}\index{conflict!vertex}
at time $t$ if $\pi_i(t) = \pi_j(t)$, and a \textbf{swapping
conflict}\index{conflict!swapping} at time $t$ if $\pi_i(t) = \pi_j(t+1)$ and
$\pi_i(t+1) = \pi_j(t)$. A plan is valid if no pair of paths has a conflict of either
kind. Following conflicts, where an agent enters the cell that another agent has just
left, are allowed, as in the classical setting.

\begin{definition}[Reservation table]
\label{def:ch08-reservation}
A \textbf{reservation table}\index{reservation table!definition} is a pair
$R = (R_V, R_E)$ of sets
\[
  R_V \subseteq V \times \N, \qquad R_E \subseteq E' \times \N,
\]
where $E'$ contains both orientations $(u, v)$ and $(v, u)$ of every edge. An entry
$(v, t) \in R_V$ means that some higher-priority agent occupies vertex $v$ at time $t$;
an entry $((u, v), t) \in R_E$ means that some higher-priority agent moves from $u$ to
$v$ between $t$ and $t+1$, so that moving from $v$ to $u$ in that step would be a swap.
Reserving a path $\pi$ with arrival time $T$ adds $(\pi(t), t)$ for
$0 \le t \le T$, adds $((\pi(t), \pi(t+1)), t)$ for every move $\pi(t) \ne \pi(t+1)$,
and marks the goal $\pi(T)$ as \textbf{parked}\index{reservation table!parked agent}
from time $T$ on, which stands for the infinite set of entries $\set{(\pi(T), t) : t \ge T}$.
\end{definition}

A space-time state $(v, t)$ is \textbf{blocked}\index{reservation table!blocked state}
by $R$ if $(v, t) \in R_V$ or $v$ is parked on at a time $\le t$. A move from $u$ to $v$
at time $t$ is blocked if $(v, t+1)$ is blocked or $((v, u), t) \in R_E$. Waiting at
$u$ is blocked if $(u, t+1)$ is blocked.

\begin{definition}[Priority order and prioritized plan]
\label{def:ch08-order}
A \textbf{priority order}\index{priority order} is a permutation $\sigma$ of
$\set{1, \dots, k}$; agent $\sigma(1)$ has the highest priority. The
\textbf{prioritized plan}\index{prioritized planning!plan} for $\sigma$ is obtained by
planning the agents in the order $\sigma(1), \sigma(2), \dots$ where agent $\sigma(j)$
receives a shortest path that is not blocked by the reservation table containing the
paths of $\sigma(1), \dots, \sigma(j-1)$. If some agent has no such path, the plan for
$\sigma$ \emph{fails}. Under the \textbf{revised rule}\index{prioritized planning!revised rule}
of \v{C}\'ap et al.~\cite{cap2015prioritized}, agent $\sigma(j)$ additionally treats
the start vertices $s_{\sigma(j+1)}, \dots, s_{\sigma(k)}$ of the lower-priority agents
as static obstacles.
\end{definition}

The definition makes precise what \emph{greedy} means here: each agent takes a path that
is optimal for itself given the agents before it, and that path is never revised. The
revised rule costs nothing in open maps and is the ingredient that makes the
completeness theorem of \cref{sec:ch08-properties} work: a lower-priority agent that
has not moved yet can never be run over.

\begin{definition}[Well-formed instance]
\label{def:ch08-wellformed}
Let $P = \set{s_1, \dots, s_k, g_1, \dots, g_k}$ be the set of all
\textbf{endpoints}\index{endpoint}. An instance is
\textbf{well-formed}\index{well-formed instance} if for every agent $i$ there is a path
in $G$ from $s_i$ to $g_i$ that visits no endpoint of $P$ other than $s_i$ and $g_i$. In
words: every agent can reach its goal without passing through the start or goal of any
other agent~\cite{cap2015prioritized}.
\end{definition}

Well-formedness is a property of the map and the endpoints alone; it does not depend on
time or on the priority order, and it is cheap to test with $k$ breadth-first searches
(\cref{ch:ch03}) on the graph with the other endpoints removed. Most warehouse layouts,
where starts and goals are parking bays off the main aisles, are well-formed by design.

% ---------------------------------------------------------------------
\section{The algorithm}
\label{sec:ch08-algorithm}

\Cref{alg:ch08-pp} is the outer loop. It takes an order, keeps one reservation table, and
calls Cooperative \astar (\cref{alg:ch08-ca}) once per agent. Cooperative \astar is
the space-time \astar of \cref{ch:ch04} with three changes: successors are filtered by
the reservation table, the goal test also requires that the goal is free for all later
times, and the search is cut off at a time horizon $T_{\max}$ so that it terminates when
no path exists.

\begin{algorithm}[htb]
\caption{Prioritized planning}
\label{alg:ch08-pp}
\KwIn{graph $G$, agents with starts $s_i$ and goals $g_i$, priority order $\sigma$, horizon $T_{\max}$, flag \emph{revised}}
\KwOut{a conflict-free plan $(\pi_1, \dots, \pi_k)$ or \emph{failure}}
\Fn{\PpPlan{$G$, $s$, $g$, $\sigma$, $T_{\max}$}}{
  $R \gets (\emptyset, \emptyset)$\tcp*{empty reservation table}
  \For{$j \gets 1$ \KwTo $k$}{
    $i \gets \sigma(j)$\label{alg:ch08-pp:pick}\;
    \lIf{revised}{$B \gets \set{s_{\sigma(m)} : m > j}$ \KwElse $B \gets \emptyset$\label{alg:ch08-pp:avoid}}
    $\pi_i \gets$ \PpStar{$G$, $s_i$, $g_i$, $R$, $B$, $T_{\max}$}\label{alg:ch08-pp:search}\;
    \If{$\pi_i = $ failure}{\KwRet{failure}\label{alg:ch08-pp:fail}\;}
    \PpReserve{$R$, $\pi_i$}\label{alg:ch08-pp:reserve}\tcp*{vertices, reverse edges, parked goal}
  }
  \KwRet{$(\pi_1, \dots, \pi_k)$}\;
}
\end{algorithm}

\begin{algorithm}[htb]
\caption{Cooperative \astar: space-time \astar against a reservation table}
\label{alg:ch08-ca}
\KwIn{graph $G$, start $s$, goal $g$, reservation table $R$, static obstacles $B$, horizon $T_{\max}$, heuristic $\hcost$}
\KwOut{a shortest unblocked path from $s$ to $g$ or \emph{failure}}
\Fn{\PpStar{$G$, $s$, $g$, $R$, $B$, $T_{\max}$}}{
  $t_g \gets$ last time at which $g$ is reserved in $R$ ($-1$ if never, $\infty$ if parked on)\label{alg:ch08-ca:tg}\;
  \lIf{$t_g = \infty$ \KwOr $(s, 0)$ is blocked}{\KwRet{failure}}
  \Open $\gets \set{(s, 0)}$; \Closed $\gets \emptyset$\;
  \While{\Open $\neq \emptyset$}{
    $(v, t) \gets \argmin_{\text{\Open}} \fcost(v, t) = t + \hcost(v)$\label{alg:ch08-ca:pop}\;
    move $(v, t)$ from \Open to \Closed\;
    \If{$v = g$ \KwAnd $t > t_g$}{\KwRet{path reconstructed from parents}\label{alg:ch08-ca:goal}\;}
    \lIf{$t \ge T_{\max}$}{\KwContinue\label{alg:ch08-ca:horizon}}
    \ForEach{$v' \in \set{v} \cup \Succ(v)$, $v' \notin B$\label{alg:ch08-ca:succ}}{
      \lIf{$(v', t+1)$ is blocked in $R$}{\KwContinue\label{alg:ch08-ca:vertex}}
      \lIf{$v' \ne v$ \KwAnd $((v', v), t) \in R_E$}{\KwContinue\label{alg:ch08-ca:edge}}
      \If{$(v', t+1) \notin$ \Open $\cup$ \Closed}{
        parent$(v', t+1) \gets (v, t)$; insert $(v', t+1)$ into \Open\label{alg:ch08-ca:insert}\;
      }
    }
  }
  \KwRet{failure}\;
}
\end{algorithm}

\paragraph{Walkthrough of \cref{alg:ch08-pp}.} Line~\ref{alg:ch08-pp:pick} picks the
next agent in priority order. Line~\ref{alg:ch08-pp:avoid} builds the set $B$ of static
obstacles: empty in Silver's original method, the starts of all lower-priority agents
under the revised rule. Line~\ref{alg:ch08-pp:search} runs a full single-agent search;
the reservation table $R$ is the only channel through which earlier agents influence
it. If the search fails (line~\ref{alg:ch08-pp:fail}) the whole procedure fails: there
is no backtracking, no re-ordering and no replanning of earlier agents. The callers of
\cref{sec:ch08-variants} add restarts with new orders around this loop.
Line~\ref{alg:ch08-pp:reserve} reserves the new path exactly as
\cref{def:ch08-reservation} says.

\paragraph{Walkthrough of \cref{alg:ch08-ca}.} Line~\ref{alg:ch08-ca:tg} computes the
\textbf{goal-stay time}\index{goal-stay check} $t_g$: the last time step at which some
higher-priority agent visits $g$. If a higher-priority agent is parked on $g$, then $g$
is reserved for ever and the search cannot succeed, so the function returns failure
immediately. The goal test in line~\ref{alg:ch08-ca:goal} accepts $(g, t)$ only if
$t > t_g$: from then on nobody else will use $g$, so the agent can stay there. Because
every action costs one, the cost so far of a state $(v, t)$ is simply $t$, which is why
the key in line~\ref{alg:ch08-ca:pop} is $t + \hcost(v)$ and why the first time a
space-time state is generated is also the cheapest: no state in \Open ever needs its
cost updated (line~\ref{alg:ch08-ca:insert}). Line~\ref{alg:ch08-ca:succ} generates the
wait action first and then the moves, skipping the static obstacles in $B$;
line~\ref{alg:ch08-ca:vertex} discards successors whose target cell is reserved, and
line~\ref{alg:ch08-ca:edge} discards moves that would swap with a reserved move. The
horizon test in line~\ref{alg:ch08-ca:horizon} stops the search from waiting for ever:
without it, a blocked agent would generate $(s, 1), (s, 2), \dots$ without end. A safe
choice of $T_{\max}$, and a trick that removes the need for it, are discussed in
\cref{sec:ch08-implementation}.

The heuristic $\hcost$ must be admissible for the static single-agent problem; the
Manhattan distance is the usual first choice on 4-connected grids
(\cref{ch:ch04}). Reservations only remove states, so a heuristic that is admissible
without reservations is admissible with them, and \cref{alg:ch08-ca} returns a path
that is shortest among all unblocked paths whose arrival time is at most $T_{\max}$.
Ties between states with equal $\fcost$ are broken towards the larger $t$, the
tie-breaking rule of \cref{ch:ch04}, and then in order of generation, so that the
search prefers to keep moving rather than to wait.

% ---------------------------------------------------------------------
\section{A worked example}
\label{sec:ch08-example}

\begin{example}[Three agents and a short corridor]
\label{ex:ch08-three}
\Cref{fig:ch08-example} shows a $3 \times 7$ grid. Cells are written $(r, c)$ with the
row $r \in \set{0, 1, 2}$ counted from the top and the column $c \in \set{0, \dots, 6}$
from the left, as in \cref{ch:ch07}. The middle row is blocked except for its three
open cells $(1,0)$, $(1,3)$ and $(1,6)$, so the map consists of two corridors joined at
both ends and in the middle. Agent~1 flies along the bottom corridor from $s_1 = (2,0)$
to $g_1 = (2,6)$, agent~2 along the top corridor from $s_2 = (0,0)$ to $g_2 = (0,6)$,
and agent~3 crosses from $s_3 = (0,3)$ down through the middle opening to
$g_3 = (2,3)$. Alone, the agents need $6$, $6$ and $2$ steps; the sum $14$ is a lower
bound on every valid plan. The instance is not well-formed: every path of agent~3 ends
on the bottom corridor, which agent~1 cannot avoid.
\end{example}

\begin{figure}[tb]
  \centering
  \inputfigure{ch08/example}
  \caption{\Cref{ex:ch08-three} in two priority orders. Small numbers are arrival times.
  (a) In the order $(1,2,3)$ the two corridor agents take their straight paths; agent~3
  reaches its goal at $t = 2$, but must step back up at $t = 3$ to let agent~1 pass
  underneath and returns at $t = 4$ (dashed). (b) In the order $(3,1,2)$ agent~3 parks
  on the bottom corridor at $t = 2$, agent~1 must take the top corridor ($10$ steps),
  and agent~2 has to dodge into $(1,3)$ at $t = 5$ to let agent~1 through before it can
  follow it to its goal ($9$ steps).}
  \label{fig:ch08-example}
\end{figure}

\paragraph{Building the table.} Take the order $(1,2,3)$. Agent~1 plans against an
empty table and gets the straight path $\pi_1(t) = (2, t)$ for $t = 0, \dots, 6$. Its
reservation adds the seven vertex entries $((2,t), t)$, the six move entries
$(((2,t), (2,t+1)), t)$ for $t = 0, \dots, 5$, and parks $(2,6)$ from $t = 6$ on. Agent~2
then plans against this table. Nothing in it touches the top corridor, so it gets
$\pi_2(t) = (0, t)$ and reserves the corresponding entries. \Cref{tab:ch08-table} lists
the table that agent~3 now faces. Only three of its entries matter for agent~3, all in
column $c = 3$: the vertices $(2,3)$ and $(0,3)$ are taken at $t = 3$, and the move
$(2,2) \to (2,3)$ at $t = 2$ forbids the reverse move at that time. The goal-stay time
of agent~3 is $t_g = 3$, the last time at which $(2,3)$ is reserved.

\begin{table}[tb]
  \centering\small
  \caption{The reservation table of \cref{ex:ch08-three} after agents~1 and~2 have been
  planned in the order $(1,2,3)$. The third row lists the entries that affect agent~3.}
  \label{tab:ch08-table}
  \begin{tabular}{@{}lL{5.2cm}L{4.6cm}l@{}}
  \toprule
  Agent & Vertex entries $R_V$ & Move entries $R_E$ & Parked \\
  \midrule
  1 & $((2,t), t)$ for $t = 0, \dots, 6$ & $(((2,t),(2,t{+}1)), t)$ for $t = 0, \dots, 5$ & $(2,6)$ from $t = 6$ \\
  2 & $((0,t), t)$ for $t = 0, \dots, 6$ & $(((0,t),(0,t{+}1)), t)$ for $t = 0, \dots, 5$ & $(0,6)$ from $t = 6$ \\
  \midrule
  seen by 3 & $((2,3), 3)$, $((0,3), 3)$ & $(((2,2),(2,3)), 2)$: no move $(2,3) \to (2,2)$ at $t = 2$ & $t_g = 3$ \\
  \bottomrule
  \end{tabular}
\end{table}

\paragraph{Running Cooperative \astar for agent~3.} \Cref{tab:ch08-trace} traces
\cref{alg:ch08-ca} for agent~3 with the Manhattan heuristic, and
\cref{fig:ch08-spacetime} draws the same search as a space-time diagram of column
$c = 3$. The first two expansions go straight down. At step~3 the search pops
$((2,3), 2)$: the goal, but too early, because $t = 2 \le t_g = 3$. Its three
successors show all three kinds of blocking at once: waiting at $(2,3)$ until $t = 3$
is a vertex conflict with agent~1, moving left to $(2,2)$ would swap with agent~1's
move $(2,2) \to (2,3)$, and only moving up to $(1,3)$ and right to $(2,4)$ survive.
Steps~4 and~5 pop the two remaining states with $\fcost = 3$ and find them blocked as
well (moving up from $(1,3)$ at $t = 2$ runs into agent~2 at $(0,3)$). Step~6 expands
$((1,3), 3)$, and step~7 pops $((2,3), 4)$ with $t = 4 > t_g$: the goal test succeeds.
The path is $(0,3), (1,3), (2,3), (1,3), (2,3)$ with cost $4$, two more than the
undisturbed $2$, and the sum of costs of the order $(1,2,3)$ is $6 + 6 + 4 = 16$.

\begin{table}[tb]
  \centering\footnotesize
  \caption{Trace of \cref{alg:ch08-ca} for agent~3 of \cref{ex:ch08-three} in the order
  $(1,2,3)$, produced by \code{code/ch08\_prioritized.py}. Generated successors are
  space-time states at time $t+1$; V marks a vertex reservation, E a swap with a
  reserved move. Ties in $\fcost$ are broken towards the larger $t$.}
  \label{tab:ch08-trace}
  \begin{tabular}{@{}rlcccll@{}}
  \toprule
  Step & Expanded $(v, t)$ & $t$ & $\hcost$ & $\fcost$ & Generated & Blocked \\
  \midrule
  1 & $((0,3), 0)$ & 0 & 2 & 2 & $(0,3)$ wait, $(1,3)$, $(0,2)$, $(0,4)$ & -- \\
  2 & $((1,3), 1)$ & 1 & 1 & 2 & $(1,3)$ wait, $(0,3)$, $(2,3)$ & -- \\
  3 & $((2,3), 2)$ & 2 & 0 & 2 & $(1,3)$, $(2,4)$ & $(2,3)$ wait: V, $(2,2)$: E \\
  4 & $((1,3), 2)$ & 2 & 1 & 3 & $(1,3)$ wait & $(0,3)$: V, $(2,3)$: V \\
  5 & $((0,3), 1)$ & 1 & 2 & 3 & $(0,3)$ wait, $(0,4)$ & $(0,2)$: V \\
  6 & $((1,3), 3)$ & 3 & 1 & 4 & $(1,3)$ wait, $(0,3)$, $(2,3)$ & -- \\
  7 & $((2,3), 4)$ & 4 & 0 & 4 & goal: $t = 4 > t_g = 3$ & \\
  \bottomrule
  \end{tabular}
\end{table}

\begin{figure}[tb]
  \centering
  \inputfigure{ch08/spacetime}
  \caption{The reservation table of \cref{tab:ch08-table} restricted to column $c = 3$,
  drawn in space-time: one row per cell, one column per time step. Shaded squares are
  entries of $R_V$ (agent~2 at $(0,3)$ and agent~1 at $(2,3)$, both at $t = 3$), the
  arrows below the diagram are agent~1's moves through $(2,3)$. The green path is the
  one agent~3 finds; the dashed red arrows are the successors that the table forbids at
  $t = 2$. From $t = 4$ on agent~3 is parked at its goal (light green).}
  \label{fig:ch08-spacetime}
\end{figure}

\paragraph{All six orders.} \Cref{tab:ch08-orders} gives the result of
\cref{alg:ch08-pp} for every permutation. Whenever agent~3 is planned last or second,
the plan costs $16$, which is also the optimal sum of costs of the instance (the code
checks this with a brute-force search over the joint space of the three agents). The
orders that start with agent~3 are a different story. Agent~3 plans alone, takes two
steps down and parks on the bottom corridor at $t = 2$ for ever. In the order
$(3,1,2)$, agent~1 must now go around through the top corridor: up to $(0,0)$, right
along row~$0$, and down through $(1,6)$, which takes $10$ steps and reserves the top
corridor at the very times agent~2 wants to use it. Agent~2 follows agent~1 to $(0,3)$,
waits there while agent~1 catches up, steps aside into $(1,3)$ at $t = 5$, lets
agent~1 pass, and follows it to $(0,6)$, arriving at $t = 9$ just after agent~1 has left
towards its own goal. The sum of costs is $10 + 9 + 2 = 21$, five more than optimal,
because a two-step agent was allowed to cut the corridor that a six-step agent needed.
In the order $(3,2,1)$ it is even worse: agent~2 parks on $(0,6)$ at $t = 6$, so the
detour through the top corridor is closed as well, and agent~1 has no path at all. The
same happens in the order $(2,3,1)$. Two of the six orders fail, one is expensive, three
are optimal, and nothing in the instance tells you in advance which is which.

\begin{table}[tb]
  \centering\small
  \caption{Sum of costs of the prioritized plan of \cref{ex:ch08-three} for every
  priority order (from \code{code/ch08\_prioritized.py}). The optimal sum of costs is
  $16$, and the lower bound from independent shortest paths is $14$.}
  \label{tab:ch08-orders}
  \begin{tabular}{@{}lcccll@{}}
  \toprule
  Order & $T_1$ & $T_2$ & $T_3$ & $\sumcost$ & What happens \\
  \midrule
  $(1,2,3)$ & 6 & 6 & 4 & 16 & agent~3 steps aside once \\
  $(1,3,2)$ & 6 & 6 & 4 & 16 & as above; agent~2 is unaffected \\
  $(2,1,3)$ & 6 & 6 & 4 & 16 & as above \\
  $(2,3,1)$ & -- & 6 & 2 & failure & both corridors closed for agent~1 \\
  $(3,1,2)$ & 10 & 9 & 2 & 21 & agent~1 detours, agent~2 dodges and waits \\
  $(3,2,1)$ & -- & 6 & 2 & failure & both corridors closed for agent~1 \\
  \bottomrule
  \end{tabular}
